Deep Reinforcement Learning (DRL) combines reinforcement learning
\cite{sutton2018rl} with function approximation and has become a standard
approach for learning control policies from interaction data. In Deep
Q-Learning (DQN)~\cite{mnih2015dqn}, a neural network is used to approximate
action-value functions, making it possible to act in environments with
continuous or high-dimensional observations. The DQN framework has been
extended with techniques such as Double Q-Learning~\cite{vanhasselt2016double}
to reduce overestimation bias and Dueling architectures~\cite{wang2016dueling}
to improve representation efficiency. Further improvements~\cite{hessel2018rainbow}
include prioritized replay and distributional RL; this work focuses on Double
and Dueling as the most architecture-relevant extensions for VQC-based Q-functions.

In parallel, Variational Quantum Circuits (VQCs), also known as parametrized
quantum circuits, have been studied as trainable models for near-term quantum
machine learning. VQCs encode classical data into a quantum state, apply
sequences of trainable unitary transformations, and extract classical outputs
via measurement. Their differentiability through parameter-shift rules~\cite{schuld2019parameter}
or backpropagation~\cite{perez2020reuploading} makes them compatible with
gradient-based optimization in standard deep learning frameworks.

Recent work has investigated replacing classical neural networks with VQCs in
reinforcement learning~\cite{chen2020vqc,jerbi2021pqc,skolik2022quantum}.
These studies suggest that VQCs can serve as Q-value approximators in DQN-style
algorithms, though the conditions under which they offer practical benefits
remain unclear, particularly under tight parameter budgets.

This work evaluates a conservative question: if the Q-network of a DQN agent is
replaced by a small VQC, and the classical baseline is constrained to a similar
number of trainable parameters, how do the two agents compare in sample
efficiency, stability, and training time on a simple control task?

The specific research questions are:

\begin{enumerate}
  \item Under a matched parameter budget of approximately 22--23 trainable
        parameters, do DQN-MLP and DQN-VQC agents show distinguishable
        performance on CartPole-v1?
  \item Do Double DQN and Dueling DQN variants alter the relative ranking
        between classical and quantum Q-functions at this budget?
  \item What is the practical simulation cost of using a VQC as a Q-function
        approximator, and can batched circuit evaluation reduce this overhead?
  \item To what extent does the representational capacity of small VQC models
        exceed what the DQN reward signal can effectively exploit, as measured
        by supervised distillation from a known-good teacher?
\end{enumerate}

The goal is not to claim quantum advantage. The experiments are performed with
a classical simulator and a small number of qubits. Instead, the contribution
is a controlled implementation and empirical comparison of a DQN-VQC agent
against a parameter-matched MLP-DQN baseline, including Double and Dueling
variants with VQC backends, a batched QNode optimization that substantially
reduces simulation overhead, a transparent report of negative or inconclusive
results, a supervised distillation experiment isolating representational
capacity from reward-based learning, and a fully reproducible codebase with
YAML-configurable experiments and comprehensive tests.

This framing is important because small simulated VQCs are not expected to
outperform classical networks in this setting. The relevant question is whether
the hybrid agent can be integrated into a standard DQN workflow and what
practical costs or instabilities appear under a controlled parameter budget.
The paper is organized as follows. Section~2 reviews background and related
work. Section~3 describes the methodology, including DQN training, VQC
architecture, and DQN variants. Section~4 presents the experimental setup.
Section~5 reports and discusses the results. Section~6 lists limitations,
and Section~7 concludes.
